\section*{Supplementary Note: Dynamic Estimation of Inhibitory Reversal Potential ($E_i$)}

The synaptic conductance reconstruction method relies on the precise knowledge of the inhibitory reversal potential ($E_i$). While theoretical values (e.g., --80~mV) are often used, the effective $E_i$ can vary significantly across preparations. To ensure assumption-free reconstructions, we implemented a data-driven geometric algorithm to estimate $E_i$ directly from the experimental data.

\subsection*{Geometric Rationale}
The relationship between total conductance ($G_{tot}(\phi)$) and the current y-intercept ($I_0(\phi)$) defines a ``wedge'' in the $G_{tot}$--$I_0$ plane. In the absence of excitatory drive ($g_e(\phi) = 0$), the current balance equation dictates a linear relationship whose slope is determined by $E_i$:
\begin{equation} \label{eq:ei_boundary}
    I_0(\phi) = - E_i G_{tot}(\phi) + I_{Li}
\end{equation}
where $I_{Li} = g_L(E_i - E_L)$ is the leak current at the inhibitory reversal potential as defined in the Methods. Geometrically, the line of pure inhibition defined by Equation~\ref{eq:ei_boundary} forms a strict upper boundary for the data trajectory in the wedge plot.

\subsection*{Tight-Fit Geometric Algorithm}
To estimate $E_i$ robustly, we use a geometric pivoting algorithm that identifies the optimal value of $E_i$ such that Equation~\ref{eq:ei_boundary} acts as a global upper bound for all observed regression points $(G_{tot}(\phi), I_{0}(\phi))$. The algorithm selects $E_i$ to minimize the area between this boundary and the data points while observing the global bounding constraint. This approach provides a physically motivated, preparation-specific estimate of $E_i$, ensuring that the subsequent decomposition into excitatory and inhibitory conductances is consistent with the global statistical limits of the recording.
